% 结论：黑体小三居中，“结  论”中间空 2 字；单独成页且限一页
\xuptheading[结论]{结\hspace{2\ccwd}论}

本文围绕校园网环境下的网络异常检测问题，设计并实现了一种基于深度学习的检测系统。首先分析了传统方法在未知攻击检测方面的不足，明确了引入深度学习方法的必要性；其次构建了以流量统计特征为基础的输入向量，采用深度自编码器学习正常流量的隐含分布，并通过注意力机制提升模型对关键特征的敏感度；最后完成了数据采集、特征工程、模型训练与可视化告警四个模块的实现，并在公开数据集上进行了对比测试。

实验结果表明，本文方法在准确率、召回率与 $F_1$ 值三项指标上均优于规则匹配与支持向量机等对比方法，单批数据检测延迟约为 18 ms，能够满足实时检测的工程需求。系统采用松耦合架构，各模块可独立升级，具备较好的可扩展性。

本文工作仍存在进一步完善的空间：一方面，训练样本主要来自公开数据集，后续可引入真实校园网流量进行迁移学习；另一方面，模型在极端流量突发场景下的误报率仍有下降空间，这将是下一步研究的重点。
